% ======================================================================
%  SECTION --- FONCTION DIGAMMA_SAVARD/CLINE (20/10/2026)
%  Fragment insere avant l'annexe.
% ======================================================================
\section{Fonction Digamma\_Savard/Cline --- 20/10/2026}
\addcontentsline{toc}{section}{Fonction Digamma\_Savard/Cline --- 20/10/2026}

\begin{center}
\begin{tikzpicture}[scale=1]
  \node[draw=spectralbleu,fill=spectralgris,rounded corners=2pt,inner sep=6pt,
        font=\small\bfseries,text width=3.5cm,align=center] (A) at (-5.25,0)
        {Observation empirique\\ des quatre possibilités\\ \textcolor{spectralrouge}{(non calculable)}};
  \node[draw=spectralor,fill=spectralor!12,rounded corners=2pt,inner sep=6pt,
        font=\small\bfseries,text width=3.5cm,align=center] (B) at (-1.75,0)
        {Réduction algébrique\\ des quatre branches\\ \textcolor{spectralbleu}{(polynôme $B_0$)}};
  \node[draw=spectralviolet,fill=spectralviolet!12,rounded corners=2pt,inner sep=6pt,
        font=\small\bfseries,text width=3.5cm,align=center] (C) at (1.75,0)
        {Criblage modulaire\\ survivant unique\\ \textcolor{spectralviolet}{(100\,\% exact)}};
  \node[draw=spectralbleu,fill=spectralbleu!14,rounded corners=2pt,inner sep=6pt,
        font=\small\bfseries,text width=3.5cm,align=center] (D) at (5.25,0)
        {Digamma calculé\\ et reconstruction\\ \textcolor{spectralbleu}{$P$}};
  \draw[-{Stealth[length=2.4mm]},spectralbleu,thick] (A)--(B);
  \draw[-{Stealth[length=2.4mm]},spectralbleu,thick] (B)--(C);
  \draw[-{Stealth[length=2.4mm]},spectralbleu,thick] (C)--(D);
  \node[font=\scriptsize,text=black!60,below=2mm of A,align=center]
        {choix de la branche\\ et rigidité de la preuve};
\end{tikzpicture}
\end{center}

\begin{boitenote}[Fil conducteur de la section]
Cette section ne réécrit pas la méthode spectrale : elle en documente
l'\emph{unique} pièce qui manquait à la forme. Le Digamma
calculé --- la correction qui transforme une somme de suite A en nombre
premier --- jusqu'alors était une \emph{valeur observée}, sélectionnée
par l'œil et la comparaison. Il devient ici le résultat d'une
\emph{fonction calculable}, dont on donne la définition, la dérivation,
la validation numérique et les limites.
\end{boitenote}

\subsection{Problématique : de l'observation empirique à la fonction calculable}

Avant cette section, la valeur du \textbf{Digamma calculé} n'était pas une
valeur \emph{calculée} : c'était une valeur \emph{repérée}. Le raisonnement
tenait en quatre temps.

\begin{enumerate}[nosep,leftmargin=*]
  \item \textbf{Quatre possibilités, et seulement quatre.} Pour un rapport
    typique $t = 1/k$, le Digamma pouvait être la somme de la suite A
    corrigée par $\pm\,t^{8}$ ou par $\pm\,t^{7}$, soit quatre
    possibilités --- la 8\textsuperscript{ième} position de la suite A
   \emph{ajoutée ou soustraite}, la 7\textsuperscript{ième} position
    \emph{ajoutée ou soustraite}. La polarité (positivité ou négativité)
    du terme était donc un \emph{choix}, non une \emph{conséquence}.
  \item \textbf{Un intervalle de calibration.} Ce choix ne pouvait être
    départagé que sur l'intervalle des quatre premiers
    $P = 29,\ 31,\ 37,\ 41$, où la valeur du Digamma était à la fois
    \emph{calculable} et \emph{raisonnable}, c'est-à-dire compatible avec
    une récurrence exhibiting une logique interne.
  \item \textbf{Une récurrence observée.} Les trois valeurs suivantes
    s'obtenaient l'une de l'autre par addition de $+4$ et par
    consommation successive des trois derniers termes de la suite A
    ($256$, $384$, $768$).
  \item \textbf{Une généralisation contestable.} Le Digamma calculé
    étant dès lors fixé, on en déduisait la reconstruction
    \[
      P \;=\; \frac{\SB - \mathrm{Digamma}_{\mathrm{calculé}}}{k^{6}},
    \]
    formule qui, appliquée à tout autre $P$, donnait « un Digamma
    calculé également déterminable » --- mais dont la valeur était
    héritée de l'observation, non \emph{démontrée}.
\end{enumerate}

\begin{boitealerte}[Le point de rupture]
Cette Circularité est aujourd'hui \textbf{levée sur le point
essentiel} : le signe et la position ne sont plus choisis, ils sont
\emph{déduits} par un argument arithmétique qui ne fait intervenir
aucune observation. En revanche --- et c'est explicité plus bas --- le
reste de la chaîne conserve un statut \emph{empirique} qui doit être lu
comme tel.
\end{boitealerte}

\subsection{Formulation algébrique du problème}

Rappelons les deux sommes fermées de la méthode spectrale (niveau~2) :
\[
  \SA(k,n) = \frac{\alpha_A(k)}{2}\,k^{n} - \delta_A(k),
  \qquad
  \SB(k,n) = \frac{\alpha_B(k)}{2}\,k^{n} - \delta_B(k),
\]
et les quatre valeurs candidates du Digamma au niveau~4 :
\[
  \mathrm{Digamma}(k,n) \;=\; \SA(k,n) \pm k^{8}
  \qquad\text{ou}\qquad
  \mathrm{Digamma}(k,n) \;=\; \SA(k,n) \pm k^{7}.
\]
La question que la section tranche est la suivante :

\begin{boiteformule}[Problème de la sélection de branche]
Étant donné $k$ et $n$, \emph{comment décide-t-on}, par un calcul et non
par une observation, que le Digamma correct est $\SA - k^{8}$ plutôt que
$\SA + k^{8}$, $\SA - k^{7}$ ou $\SA + k^{7}$~?
\end{boiteformule}

Le point crucial est que l'on \emph{ne demande pas} « le candidat est-il
premier~? » (question qui ramènerait la circularité) mais « le candidat
est-il \emph{éliminé} par un petit premier~? ». C'est une question de
\emph{division}, non de \emph{primalité} : elle admet une réponse
modulaire, donc calculable.
\subsection{Définition formelle de la fonction \textsc{Digamma\_Savard}}

\subsubsection{Réduction algébrique : le polynôme $B_0(k)$}

L'observation décisive est que la reconstruction
\[
  P \;=\; \frac{\SB - \mathrm{Digamma}}{k^{6}}
\]
se simplifie : $\SB - \SA$ est divisible par $k^{6}$ pour tout $k$, et le
quotient est un polynôme en $k$ à cinq termes. En effet, au niveau
$n = 10$,
\begin{align*}
  \SA(k,10) &= k^{10}+k^{9}+k^{6}+k^{5}+k^{4}+k^{3}+k^{2}+k,\\
  \SB(k,10) &= k^{11}+k^{10}+k^{7}+k^{5}+k^{4}+k^{3}+k^{2}+k,\\
  \SB(k,10)-\SA(k,10) &= k^{11}-k^{9}+k^{7}-k^{6}
  \;=\; k^{6}\bigl(k^{5}-k^{3}+k-1\bigr).
\end{align*}

\begin{boiteformule}[Polynôme de base $B_0$]
\[
  \boxed{\;B_0(k) \;=\; \frac{\SB(k,10)-\SA(k,10)}{k^{6}}
  \;=\; k^{5}-k^{3}+k-1\;}
\]
\end{boiteformule}

Le Digamma étant $\SA \pm k^{8}$ ou $\SA \pm k^{7}$, on obtient, pour
chacune des quatre branches, un candidat \emph{polynomial} en $k$ :
\begin{align*}
  P_{A(8)-}(k) &= B_0(k) + k^{2}, &
  P_{A(8)+}(k) &= B_0(k) - k^{2},\\
  P_{A(7)+}(k) &= B_0(k) - k, &
  P_{A(7)-}(k) &= B_0(k) + k.
\end{align*}

\begin{boitenote}[Pourquoi la polarité de la suite A commande le signe]
Le signe du Digamma et le signe de la correction \emph{retombent} sur la
polarité de la position dans la suite~A : la position~8 (terme
$k^{8}$) donne une correction en $k^{2}$ \emph{positive} quand elle est
soustraite, la position~7 (terme $k^{7}$) donne une correction en $k$
\emph{positive} quand elle est \emph{ajoutée}. Ce renversement n'est pas
conventionnel : il découle mécaniquement de
$\SB-\SA = k^{6}B_0$, où le terme de rang~6 est celui qui a été
\emph{retiré} à la suite~B par le saut Zêta.
\end{boitenote}

\subsubsection{Table de correspondance : branche, signe et terme}

\begin{center}
\small
\begin{tabular}{@{}l c c r@{}}
\toprule
\rowcolor{spectralgris}
Branche & Position dans la suite A & Terme correctif & Candidat polynomial \\
\midrule
$A(8)-$ & 8\textsuperscript{ième} & $-k^{8}$ & $B_0+k^{2}$ \\
$A(8)+$ & 8\textsuperscript{ième} & $+k^{8}$ & $B_0-k^{2}$ \\
$A(7)+$ & 7\textsuperscript{ième} & $+k^{7}$ & $B_0-k$ \\
$A(7)-$ & 7\textsuperscript{ième} & $-k^{7}$ & $B_0+k$ \\
\bottomrule
\end{tabular}
\end{center}

Cette table est la \emph{vraie} table de correspondance : elle a été
revérifiée numériquement sur $676\,191$ ancres, après qu'une version
antérieure (issue du script de criblage originel) eut inversé les
lignes $A(7)+$ et $A(7)-$. L'inversion était invisible au regard mais
destructrice : elle inversait le signe de la correction Digamma pour
une branche sur quatre.

\subsubsection{Les quatre ancres de référence : $P = 29,\ 31,\ 37,\ 41$}

Pour le rapport typique $t = 1/2$ (soit $k = 2$), les suites à dix termes
valent
\[
  \SA(2,10) = 2+4+8+16+32+64+128+256+384+768 = 1\,662,
\]
\[
  \SB(2,10) = 2+4+8+16+32+128+256+512+768+1536 = 3\,262,
\]
et la 6\textsuperscript{ième} position de la suite A (le « Zêta ») vaut
$2^{6} = 64$.

\begin{center}
\small
\begin{tabular}{@{}c r r r r r r@{}}
\toprule
\rowcolor{spectralgris}
$n$ & $P$ & $\SA(2,n)$ & $\SB(2,n)$ & $\mathrm{Digamma}$ & $\SA+\mathrm{Digamma}$ & Branche \\
\midrule
10 & 29  & 1\,662   & 3\,262   & $-256$   & 1\,406   & $A(8)-$ \\
11 & 31  & 3\,326   & 6\,590   & $+1\,280$ & 4\,606   & $A(8)+$ \\
12 & 37  & 6\,654   & 13\,246  & $+4\,224$ & 10\,878  & $A(8)+$ \\
13 & 41  & 13\,310  & 26\,558  & $+10\,624$ & 23\,934  & $A(8)+$ \\
\bottomrule
\end{tabular}
\end{center}

\begin{boiteformule}[Vérification de la reconstruction sur les quatre ancres]
\[
  \frac{\SB - \mathrm{Digamma}_{\mathrm{calculé}}}{64}
  \;=\;
  \frac{3262-1406}{64}=29,\quad
  \frac{6590-4606}{64}=31,\quad
  \frac{13246-10878}{64}=37,\quad
  \frac{26558-23934}{64}=41.
\]
\end{boiteformule}

\begin{boitenote}[Point d'exactitude relevé à la relecture]
La valeur du Digamma pour $n = 13$ vaut $10\,624$ et donne
$\SA + \mathrm{Digamma} = 23\,934$. La reconstruction
$(26\,558-23\,934)/64 = 41$ est donc exacte. Une saisie antérieure
portait $23\,936$ (soit $10\,626$), qui ne redonne pas $41$ mais
$40{,}96875$ : l'écart est de $2$ sur le Digamma et faute de
reproduction du nombre premier. La valeur $10\,624$ est la bonne.
\end{boitenote}

\subsubsection{La récurrence empirique observée}

Les quatre valeurs du Digamma se lisent aussi comme des décompositions
sur les trois derniers termes de la suite A de référence
($256$, $384$, $768$), avec des coefficients formant une progression
de pas $4$ ($1,\ 5,\ 9,\ 13$) :
\begin{align*}
  256     &= 1\times256,\\
  1\,280  &= 5\times256,\\
  4\,224  &= 9\times256 + 5\times384,\\
  10\,624 &= 13\times256 + 9\times384 + 5\times768.
\end{align*}
Cette décomposition canonique s'écrit
\begin{equation*}
  \mathrm{Digamma}(2,n) \;=\; \sum_{j=0}^{\max(0,\,n-11)}
  \bigl(4\,(n-10-j)+1\bigr)\,b_j,
  \qquad
  b_0 = 256,\; b_1 = 384,\; b_2 = 768,\; b_{j+1} = 2\,b_j .
\end{equation*}

\begin{boitealerte}[Statut exact de cette récurrence]
La formule ci-dessus est \emph{exacte} pour $n = 10, 11, 12, 13$ --- les
quatre valeurs de la table de référence. \textbf{Elle devient fausse
dès $n = 14$~} : la règle canonique donne $23\,936$ là où la valeur
exacte est $23\,808$ (écart $+128$), puis $+896$ en $n = 15$,
$+2\,816$ en $n = 16$, $+6\,784$ en $n = 17$ et $+14\,592$ en $n = 18$.
Elle est donc \emph{une description de la table}, \emph{non} une
loi générale. La fonction qui suit ne s'appuie pas sur elle.
\end{boitealerte}
\subsection{La règle de criblage modulaire --- formalisation de l'empirique}

\subsubsection{Le principe : éliminer par divisibilité, non par primalité}

Les quatre candidats $P_{A(8)-}$, $P_{A(8)+}$, $P_{A(7)+}$, $P_{A(7)-}$ sont des
\emph{polynômes} en $k$. Il en résulte une propriété décisive : pour
tout premier $p$,
\[
  c_j(k) \equiv 0 \!\!\pmod p
  \quad\Longleftrightarrow\quad
  c_j(k \bmod p) \equiv 0 \!\!\pmod p .
\]
Autrement dit, la divisibilité d'un candidat par $p$ ne dépend que de la
classe de $k$ modulo $p$. L'élimination d'un candidat par un petit
premier est donc une opération \emph{modulaire}, entièrement calculable
et \emph{indépendante} de toute conjecture sur les nombres premiers.

\begin{boiteformule}[Définition --- élimination modulaire]
Soit $\mathcal{P} = \{3,5,7,11,13,17,19,23\}$ un ensemble de petits
premiers. Le candidat $c_j(k)$ est dit \textbf{éliminé par criblage} s'il
existe $p \in \mathcal{P}$ tel que $p \mid c_j(k)$ ; il est dit
\textbf{survivant} si aucun $p \in \mathcal{P}$ ne le divise.
\end{boiteformule}

Puisque chaque $c_j$ est un polynôme, l'ensemble des branches éliminées
est une fonction de $k$ modulo
\[
  M \;=\; \prod_{p \in \mathcal{P}} p \;=\; 111\,546\,435 .
\]

\subsubsection{La règle de sélection}

\begin{boiteformule}[Théorème empirique --- règle de branche]
Étant donné $k$, on forme les quatre candidats
\[
  c_0 = B_0+k^{2},\qquad c_1 = B_0-k^{2},\qquad
  c_2 = B_0-k,\qquad c_3 = B_0+k .
\]
On élimine tout candidat divisible par au moins un premier de
$\mathcal{P} = \{3,5,7,11,13,17,19,23\}$. \textbf{Si le criblage laisse un
seul survivant, ce survivant est la branche correcte.}
\end{boiteformule}

\begin{boitenote}[Pourquoi l'élimination est sûre]
La branche correcte étant, par construction, celle qui donne le
nombre premier $P$, son candidat n'est divisible par aucun premier
strictement inférieur à $P$ --- et en particulier par aucun des huit
premiers de $\mathcal{P}$. \emph{Le criblage ne peut donc jamais
éliminer la vraie branche.} C'est la garantie d'\emph{intégrité} de la
règle, et elle est structurelle, non statistique.
\end{boitenote}

\begin{boitealerte}[Ce que la règle ne fait pas]
La règle \emph{ne remplace pas} un test de primalité. Elle indique
\emph{quelle} branche est compatible avec la primalité, pas que le
candidat retenu est premier. Symétriquement, elle ne prouve pas que les
trois candidats éliminés sont composés : ils le sont bien (ils ont un
diviseur trouvé), mais \emph{pour la mauvaise raison} --- par un petit
premier --- ce qui est une information \emph{plus faible} que la
composante complète.
\end{boitealerte}

\subsubsection{La fonction \textsc{Digamma\_Savard} proprement dite}

\begin{boiteformule}[Définition --- fonction \textsc{Digamma\_Savard/Cline}]
Pour $k \ge 2$ et $n$ entier positif, on pose
\[
  \mathrm{Digamma}_{\mathrm{S/C}}(k) \;=\;
  \begin{cases}
    \SA(k,n) - k^{8}, & \text{si le survivant unique est } A(8)-,\\
    \SA(k,n) + k^{8}, & \text{si le survivant unique est } A(8)+,\\
    \SA(k,n) + k^{7}, & \text{si le survivant unique est } A(7)+,\\
    \SA(k,n) - k^{7}, & \text{si le survivant unique est } A(7)-,
  \end{cases}
\]
et l'on pose $\mathrm{Digamma}_{\mathrm{S/C}}(k) = \text{indéterminé}$
lorsque le criblage laisse deux survivants ou davantage.
\end{boiteformule}

\begin{boitenote}[Caractère de la fonction : une fonction \emph{définie} à
totaux, \emph{partiellement} calculable]
La fonction ci-dessus est \emph{totale} en tant qu'application : pour
tout $k$, la règle de criblage produit un verdict. Elle est cependant
\emph{partiellement calculable} au sens où, sur une fraction non
négligeable des $k$, elle renvoie « indéterminé » : c'est le prix de
l'exigence d'\emph{exactitude} (on ne devine jamais) devant celle de
\emph{complétude} (on couvre tout).
\end{boitenote}
\subsection{Validations expérimentales}

Toutes les validations reposent sur le même corpus : les
$676\,191$ ancres uniques extraites d'une reconstruction portant sur
$4$ millions de rapports, et qui associent à chaque $k$ la branche
effective et la valeur correspondante.

\subsubsection{Validation principale : exactitude et intégrité}

\begin{center}
\small
\begin{tabular}{@{}l r@{}}
\toprule
\rowcolor{spectralgris}
Mesure & Résultat \\
\midrule
Ancres analysées & 676\,191 \\
Écarts au modèle polynomial $B_0$ & 0 \\
Violations du criblage (0 attendues) & \textbf{0} \\
Ancres à survivant unique & 71\,401 \quad (10{,}56\,\%) \\
Exactitude sur ces ancres & 71\,401/71\,401 \;=\; \textbf{100{,}000000\,\%} \\
Module de criblage $M$ & 111\,546\,435 \\
\bottomrule
\end{tabular}
\end{center}

\begin{boiteformule}[Résultat de la validation]
Sur l'ensemble du corpus, la règle de branche ne présente
\textbf{aucune violation} et atteint une \textbf{exactitude de
100\,\%} sur les $71\,401$ cas où elle se prononce. Le taux de
couverture --- soit la proportion de $k$ pour lesquels la règle désigne
une branche unique --- est de $10{,}56\,\%$.
\end{boiteformule}

\subsubsection{Contrôle anti-circularité}

Le risque d'une telle règle est d'être circulaire : si elle ne
distinguait rien, une sélection \emph{arbitraire} fonctionnerait aussi.
Le test consiste donc à \emph{permuter} les quatre signes et à mesurer
la pureté obtenue.

\begin{center}
\small
\begin{tabular}{@{}l c r@{}}
\toprule
\rowcolor{spectralgris}
Essai & Permutation des signes & Pureté \\
\midrule
Référence (règle true) & $[0,1,2,3]$ & \textbf{100{,}000000\,\%} \\
Essai 0 & $[3,1,0,2]$ & 8{,}627330\,\% \\
Essai 1 & $[1,2,3,0]$ & 0{,}000000\,\% \\
Essai 2 & $[1,0,3,2]$ & 0{,}000000\,\% \\
\bottomrule
\end{tabular}
\end{center}

\begin{boitenote}[Lecture du contrôle]
Sous permutation aléatoire, la pureté s'effondre à $0$--$8{,}6\,\%$,
très en deçà du taux de base de $28{,}84\,\%$ correspondant à la branche
majoritaire. La règle ne peut donc pas être « bonne » par hasard :
l'information qu'elle exploite est \emph{spécifique} à l'association
entre la polarité de la position dans la suite A et le signe de la
correction Digamma. C'est le test le plus important de la section.
\end{boitenote}

Un second contrôle porte sur le \emph{hors-échantillon} : appliqué à
$20\,000$ valeurs de $k$ arbitraires, hors du corpus d'ancres, la règle
désigne un survivant unique pour $4\,088$ d'entre elles, et le
candidat désigné est impair dans les $4\,088$ cas (rapport typique
$1/2$ : le Digamma et $k^{6}$ sont pairs, donc $P$ doit l'être aussi
après soustraction --- la parité est un test de cohérence immédiat).

\subsubsection{Convergence du criblage}

\begin{center}
\small
\begin{tabular}{@{}r r r r r r@{}}
\toprule
\rowcolor{spectralgris}
\# premiers & Couverture & Exactitude & Intégrité & Ambigüité & Temps \\
\midrule
9   & 10{,}56\,\% & 100{,}0000\,\% & 100{,}00\,\% & 89{,}44\,\% & 2\,\text{s} \\
25  & 22{,}69\,\% & 100{,}0000\,\% & 100{,}00\,\% & 77{,}31\,\% & 4\,\text{s} \\
168 & 42{,}39\,\% & 100{,}0000\,\% & 100{,}00\,\% & 57{,}61\,\% & 19\,\text{s} \\
430 & 49{,}65\,\% & 100{,}0000\,\% & 100{,}00\,\% & 50{,}35\,\% & 41\,\text{s} \\
\bottomrule
\end{tabular}
\end{center}

\begin{boitenote}[Ce que montre la convergence]
À mesure que l'on ajoute des petits premiers au crible, la
\emph{couverture} croît de façon monotone ($10{,}56\,\%$ à $49{,}65\,\%$),
tandis que l'\emph{exactitude} et l'\emph{intégrité} restent
parfaitement constantes à $100\,\%$. C'est le comportement attendu et
souhaitable : le criblage élimine de plus en plus de branches, mais ne
sacrifie jamais la rigueur. La couverture pourrait tendre vers $100\,\%$
en augmentant le nombre de premiers \emph{au prix} d'un coût de calcul
qui croît brutalement : c'est le compromis central de la règle.
\end{boitenote}

\subsection{Cas de référence : $P = 29,\ 31,\ 37,\ 41$}

Le tableau qui suit applique la procédure complète --- criblage, choix
de branche, calcul du Digamma, reconstruction --- aux quatre nombres
premiers d'ancrage, pour le rapport typique $1/2$.

\begin{center}
\small
\begin{tabular}{@{}r r l r r r@{}}
\toprule
\rowcolor{spectralgris}
$k$ & $\mathrm{Digamma}$ & Branche retenue & $P$ reconstruit & $P$ attendu & Verdict \\
\midrule
2 & $-256$     & $A(8)-$ & 29 & 29 & conforme \\
2 & $+1\,280$  & $A(8)+$ & 31 & 31 & conforme \\
2 & $+4\,224$  & $A(8)+$ & 37 & 37 & conforme \\
2 & $+10\,624$ & $A(8)+$ & 41 & 41 & conforme \\
\bottomrule
\end{tabular}
\end{center}

\begin{boitealerte}[Mise en garde sur la lecture de ce tableau]
Ces quatre lignes \emph{valident la procédure}, elles ne la
\emph{démontrent pas} : les quatre cas appartiennent à la table
d'observation qui a servi à formuler la règle. Leur valeur est
méthodologique --- ils prouvent que la formalisation reproduit
exactement le comportement empirique qu'elle remplace --- et non
probatoire au sens d'un test sur données inédites. La preuve
opératoire est le test anti-circularité et l'exactitude à
$100\,\%$ sur $71\,401$ ancres \emph{hors table d'observation}.
\end{boitealerte}
\subsection{Limites et statut épistémologique}

\begin{boitealerte}[Limite~1 --- la couverture n'est pas totale]
La règle se prononce sur $10{,}56\,\%$ des $k$ du corpus. Sur les
$89{,}44\,\%$ restants, le criblage laisse au moins deux survivants :
la fonction renvoie alors « indéterminé ». Le résidu est
\emph{intrinsèque} et non un défaut d'implémentation : sur un
$k$ donné, le nombre de branches survivantes ne dépend que de
$k \bmod M$, et aucune extension finie du crible ne peut faire
disparaître toutes les classes ambguës sans coût prohibitif.
\end{boitealerte}

\begin{boitealerte}[Limite~2 --- la règle n'est pas un test de primalité]
Elle indique \emph{quelle} branche est compatible, non que le candidat
retenu est premier. Pour lever l'ambiguïté résiduelle, un test de
primalité reste nécessaire --- c'est-à-dire un aller-retour vers la
primalité, que la formalisation n'a pas supprimé mais simplement
\emph{repoussé} à un second étage, plus étroit et mieux circonscrit.
\end{boitealerte}

\begin{boitealerte}[Limite~3 --- la partie algébrique conserve un statut
empirique]
La réduction $\SB-\SA = k^{6}B_0$ et la formule fermée de niveau~2 sont
établies et vérifiées. En revanche la \emph{table de correspondance}
entre polarité de la position et signe de la correction, ainsi que
l'ensemble des quatre candidats, sont le fruit d'une exploration
numérique : ils sont \emph{corrects} au sens où ils reproduisent
$676\,191$ ancres sans un seul écart, mais ils ne sont pas
\emph{démontrés} au sens d'un théorème. La section franchit la
frontière entre «~observé~» et «~calculé~» ; elle ne franchit pas celle
entre «~calculé~» et «~prouvé~».
\end{boitealerte}

\begin{boitealerte}[Limite~4 --- la récurrence de la table ne se généralise pas]
Comme établi plus haut, la formule canonique de décomposition du
Digamma est exacte pour $n = 10, 11, 12, 13$ et fausse dès $n = 14$.
Elle est conservée dans le présent document comme \emph{description
de la table d'observation}, à titre d'outil de calcul et de
contrôle, et non comme une loi.
\end{boitealerte}

\begin{boitenote}[Position dans l'architecture générale]
La fonction \textsc{Digamma\_Savard/Cline} intervient exactement là où la
méthode spectrale attendait une information manquante : au niveau~4
(mécanisme d'ancrage), où les quatre essais Digamma étaient jusqu'alors
lancés \emph{simultanément} et sans critère de choix. Elle ne modifie
ni les niveaux~1 à~3 (constantes, formules fermées, construction
terme à terme) ni le niveau~5 (reconstruction pour $n \neq 10$), et ne
modifie pas non plus la formule de reconstruction
\[
  P \;=\; \frac{\SB - \mathrm{Digamma}}{k^{6}},
\]
dont elle ne fournit que le \emph{contenu} de la correction.
\end{boitenote}


\subsection{Implémentation de référence}

\begin{boitenote}[Pseudo-code de la fonction]
\begin{Verbatim}[frame=single,fontsize=\footnotesize]
DIGAMMA_SAVARD(k, n) :
    S_A <- somme_suite_A(k, n)      # arithmétique entière
    S_B <- somme_suite_B(k, n)
    B0  <- k**5 - k**3 + k - 1
    C   <- [B0+k^2, B0-k^2, B0-k, B0+k]  # A(8)-, A(8)+, A(7)+, A(7)-

    survivants <- []
    pour j de 0 a 3 :
        si tous(C[j] mod p != 0 pour p dans [3,5,7,11,13,17,19,23]) :
            ajouter j a survivants

    si longueur(survivants) == 1 :
        j <- survivants[0]
        signe   <- -1 si j dans {0,3} sinon +1
        exposant <- 8 si j dans {0,1} sinon 7
        Digamma <- S_A + signe * k**exposant
        retourner (Digamma, (S_B - Digamma) / k**6)
    sinon :
        retourner (INDETERMINE, branches survivantes)
\end{Verbatim}
\end{boitenote}

\begin{boitenote}[Précautions d'implémentation]
\begin{itemize}[nosep]
  \item \textbf{Arithmétique entière obligatoire.} Pour $k \ge 50$, les
    valeurs de $\SA$ et $\SB$ dépassent $10^{15}$~: l'arithmétique
    flottante introduit des erreurs d'arrondi qui se propagent au
    test de divisibilité et faussent le verdict de criblage.
  \item \textbf{Le test de divisibilité précède toute hypothèse de
    primalité.} La fonction ne doit jamais court-circuiter sur un
    «~candidat préféré~»~: c'est précisément ce court-circuit qui
    réintroduirait la circularité.
  \item \textbf{Le cas indéterminé doit être rendu visible.} En cas
    d'ambiguïté, l'agent doit signaler \textsc{C\_NON\_CERTIFIÉ} et
    attendre une validation externe, conformément au contrat HOL,
    plutôt que de trancher silencieusement.
\end{itemize}
\end{boitenote}


\subsection{Conclusion}

La correction Digamma, longtemps traitée comme une valeur observée,
devient ici le résultat d'une \emph{fonction calculable} : la
\emph{fonction \textsc{Digamma\_Savard/Cline}}. Sa structure est la
suivante.

\begin{enumerate}[nosep,leftmargin=*]
  \item \textbf{Une réduction algébrique.} Les quatre candidats
    possibles se ramènent à quatre polynômes
    $B_0 \pm k^{2}$ et $B_0 \pm k$, où
    $B_0(k) = k^{5}-k^{3}+k-1$.
  \item \textbf{Une règle de sélection par criblage modulaire.} La
    polarité de la position dans la suite A commande le signe de la
    correction ; ce signe est \emph{déduit} par élimination
    modulaire, jamais \emph{choisi}.
  \item \textbf{Une validation rigoureuse.} Zéro violation
    d'intégrité sur $676\,191$ ancres~; exactitude de $100\,\%$ sur
    les $71\,401$ ancres à survivant unique~; effondrement de la
    pureté à $0$--$8{,}6\,\%$ sous permutation des signes~; stabilité
    de l'exactitude à $100\,\%$ lorsque la couverture passe de
    $10{,}56\,\%$ à $49{,}65\,\%$.
  \item \textbf{Un statut honnête.} La règle est exacte là où elle se
    prononce et le dit quand elle ne peut pas se prononcer~; elle ne
    prétend ni tester la primalité, ni démontrer la table de
    correspondance.
\end{enumerate}

% Ligne retirée: ancienne environnement boiteformule
\begin{boiteformule}[Énoncé final]
Pour tout $k \ge 2$, la fonction \textsc{Digamma\_Savard/Cline} retourne
une correction calculée, non observée, dont le signe est déterminé par
la seule arithmétique modulaire~; et elle retourne «~indéterminé~»
plutôt que de trancher au hasard. \emph{La fonction ne prétend pas
résoudre la primalité~: elle déplace le problème du \emph{choix} vers
la \emph{contamination}.}
\end{boiteformule}
% Contenu sans environnement boiteformule
Pour tout $k \ge 2$, la fonction \textsc{Digamma_Savard/Cline} retourne
une correction calculée, non observée, dont le signe est déterminé par
la seule arithmétique modulaire~; et elle retourne «~indéterminé~»
plutôt que de trancher au hasard. \emph{La fonction ne prétend pas
résoudre la primalité~: elle déplace le problème du \emph{choix} vers
la \emph{contamination}}.
